\section{The Degrees of Knowledge}

\begin{table}[h]\centering\small
\begin{tabular}{cccc}\toprule
\textbf{Degree} & ~ & \textbf{Object} & \textbf{Form} \\\midrule
Doxa & Opinion & Phenomena & Concept \\ Dianoia & Rational Knowledge & Thought & Idea \\ Episteme & Intuitive Knowledge & Noumena & Ideal \\\bottomrule
\end{tabular}
\caption{Degrees of knowledge}
\end{table}

Tradition teaches that there are three degrees of knowledge:

\begin{enumerate}

\item \textbf{Doxa} (opinion): our knowledge of the sensory or phenomenal world
\item \textbf{Dianoia} (rational knowledge): knowledge based on thinking and logical deductions
\item \textbf{Episteme} (intuition, gnosis, wisdom): direct knowledge of the essence of things
\end{enumerate}


Gnosis remains largely ignored or misunderstood in our day. Voegelin, for example, designated Gnostics as those who believed in a received knowledge, that was really on the level of ideas, or dianoia. In popular parlance, intuition refers to an unexpected insight into the phenomenal world. However, in Tradition, gnosis is direct knowledge of the noumenal that is usually acquired through intense spiritual practices.

An example is the awareness and experience of oneself as a person, that is, as a center of consciousness and will. It sounds simple, but in actual practice there is much more to it. It takes efforts and exercises to become aware of one's real I\footnote{\url{http://www.gornahoor.net/?p=57}} or true Will.

So much so, in fact, that the obvious is not noticed by those who claim to be the most intelligent and educated. For example, in an earlier post\footnote{\url{http://www.gornahoor.net/?p=210}}, we showed the example of someone who denies his will on the grounds that it defies the ``laws of physics''. I could point to other examples on the Internet of intelligent people who deny the validity of their own consciousness and subjectivity, believing it to be no more than electrochemical processes in the brain. This puts dianoia as the highest form of knowledge.

It is episteme alone that provides certain and absolute knowledge. It intuits the ideal form, which is non-sensory and transcendent to the manifested world. At the level of dianoia, the ideal becomes idea, or the object of thought.

Thinking is limited to manipulation of ideas and can never reach certainty. The theories of science are always subject to future revision and deeper que stions have been debated ad infinitum, from the past and again into the future. Thinking is always about essences and can never penetrate the mystery of existence, the individual, the unique, the non-repeatable, the miraculous.

Doxa, or sensory awareness of the natural, or phenomenal, world is enthralled and deceived by the multiplicity of manifestations. Doxa recognizes concepts through a process of abstraction that extracts the common characteristics. Ideal, idea and concept correspond to the three degrees of knowledge.

Faith is related to gnosis. It does not refer to beliefs in sensory things or a body of ideas (the two lower levels of knowledge). True faith is belief in ideals, or noumenal reality. It involves the entire personality, especially the Will. Meditating on the truths of Faith can lead to authentic Gnosis.


\flrightit{Posted on 2010-04-25 by Cologero}

\begin{center}* * *\end{center}

\begin{footnotesize}
\begin{sffamily}
\texttt{Comprehensor on 2010-04-26 at 09:12 said:}

While I reject the idea that free will defies the laws of physics, to be united with the divine is to surrender one's will to God in order that he may do God's Will. To put one's own will over and against God's would be a promethean rationalism. Thus to do one's ``true will'' would be to do God's Will.

\hfill

\texttt{Cologero on 2010-04-26 at 22:13 said:}

``We are all attached to the throne of the Supreme Being by a supple chain that binds but does not enslave us. What is admirable in the universal order of things is the action of free beings under the divine hand. Freely enslaved, they function at once by will and by necessity; they really do as they wish, but without being able to upset the general plan''.

\textasciitilde{} Joseph de Maistre, ``Considerations on France''


So our will is either free (in union with the Supreme Being) or determined. Remove the religious language, and this is precisely what Evola claims.

\hfill

\texttt{Comprehensor on 2010-04-27 at 08:46 said:}

What does Joseph de Maistre mean by ``we are all... ''?

Maistre seems to suggest not an either/or but an implication of both predermination and will. Actions of necessity do not override God's plan, for God is always superior to man, since God is all-knowing and everything is predetermined by him. However God allows man a bit of space to do as he will, to choose the light or the dark.

On the other hand as Guenon and Schuon have stated man is passive in the face of the Supreme Principle and he surrenders his own free will in order to become united with the divine. This is what I was getting at in my first comment.

\end{sffamily}
\end{footnotesize}

